\documentclass[11pt,twocolumn]{article}
\usepackage[margin=1in]{geometry}
\usepackage{amsmath,amssymb,amsthm}
\usepackage{graphicx}
\usepackage{booktabs}
\usepackage{hyperref}
\usepackage{xcolor}
\usepackage{algorithm}
\usepackage{algorithmic}
\usepackage{multirow}
\usepackage{caption}
\usepackage{subcaption}
\usepackage[numbers]{natbib}
\usepackage{url}

\hypersetup{colorlinks=true, linkcolor=blue!60!black, citecolor=green!50!black, urlcolor=blue!70!black}

\newcommand{\EN}{E_N}
\newcommand{\ETN}{E_{2N}}
\newcommand{\MTN}{M_{2N}}
\newcommand{\vfinal}{v_{\text{final}}}
\newcommand{\MRE}{\text{MRE}}

\title{\textbf{Two-Phase PageRank with Travel-Time Self-Loops\\ for Urban Traffic Distribution Estimation\\ and Event-Driven Anomaly Detection}}

\author{
  Research Team\\
  \textit{University of Utah}\\
  Salt Lake City, Utah, USA
}

\date{}

\begin{document}
\maketitle

% ================================================================
\begin{abstract}
We present a modified PageRank algorithm for estimating link-level traffic distributions on urban road networks using sparse GPS trajectory data.
The model introduces two key innovations: (1)~\emph{travel-time self-loops} that weight each link's mass retention proportionally to its traversal time $\mu \cdot \ell_i / s_i$, and (2)~a \emph{two-phase block structure} separating active driving from local dispersal.
Applied to the Salt Lake City road network (99,716 directed links, 467,911 GPS trajectories over 30~days), the model achieves a Top-100 Mean Relative Error (MRE) of 0.0312 and Pearson correlation $r = 0.997$ against smoothed ground truth.

We further introduce \emph{demand-weighted transitions}, a single parameter $\gamma$ that encodes average traffic demand into the transition matrix, enabling the damping factor to be raised from 0.80 to 0.96 (reducing ground-truth dependence from 20\% to 4\%) while \emph{improving} overall accuracy.
This is validated through temporal, day-of-week, and 5-fold cross-validation splits.

For anomaly detection, we demonstrate that the dwell-time parameter $\mu$ serves as an effective per-link control knob: reducing $\mu$ on disrupted links from 20 to 0.01 drops the crash-corridor prediction error from 179\% to 11.5\%, and game-day stadium-area error from 293\% to 7.8\%---all without using any event-day ground truth.
Cross-validated Tuesday baseline analysis confirms $\mu = 20$ is optimal on normal days ($r = 0.9999$ between data gap and model error), establishing that deviations from this default reliably signal real-world disruptions.
\end{abstract}

% ================================================================
\section{Introduction}
\label{sec:intro}

Accurate estimation of link-level traffic distributions is fundamental to transportation planning, real-time route guidance, and incident detection~\cite{wardrop1952road, jiang2009citywide}.
Traditional approaches rely on fixed sensor networks (loop detectors, cameras), which provide high temporal resolution but limited spatial coverage.
The proliferation of GPS-equipped mobile devices has created an alternative data source: sparse trajectory observations that cover the entire network but observe each link infrequently~\cite{herrera2010evaluation, zheng2014urban}.

The core challenge is converting sparse, noisy GPS traversal counts into a complete traffic distribution over all links.
In a typical 15-minute observation window, only 1--3\% of links in a metropolitan network have any GPS observations~\cite{horvitz2012prediction}.
Directly using these counts produces a distribution dominated by zeros, with no information about unobserved links.

We address this through a three-stage pipeline:
\begin{enumerate}
    \item \textbf{Graph-diffusion smoothing} propagates observed counts to neighbouring links using the road adjacency structure, filling gaps while preserving the spatial correlation of traffic patterns.
    \item \textbf{Two-Phase PageRank with travel-time self-loops} computes a stationary distribution that balances observed data (teleportation) against network topology (transition matrix), producing physically realistic traffic estimates.
    \item \textbf{Demand-weighted transitions} encode average traffic demand into the matrix, enabling high damping factors ($d > 0.90$) that reduce dependence on per-timeframe ground truth.
\end{enumerate}

The remainder of this paper is organised as follows.
Section~\ref{sec:data} describes the data pipeline.
Section~\ref{sec:model} presents the two-phase model.
Section~\ref{sec:demand} introduces demand-weighted transitions.
Section~\ref{sec:results} reports accuracy results.
Section~\ref{sec:events} demonstrates event detection and prediction.
Section~\ref{sec:mu} analyses the role of~$\mu$.
Section~\ref{sec:conclusion} concludes.


% ================================================================
\section{Data Pipeline}
\label{sec:data}

\subsection{Study Area and Network}

The study area is the greater Salt Lake City metropolitan region, represented as a directed graph $G = (V, E)$ with $|E| = N = 99{,}716$ directed links and average out-degree 2.87.
The network is derived from OpenStreetMap via SUMO, with each link carrying attributes: length~$\ell_i$ (metres), free-flow speed~$s_i$ (m/s), and lane count.
Link lengths range from 0.59\,m to 11,577\,m (mean 142\,m), and speeds from 2.2 to 33.5\,m/s.

\subsection{GPS Trajectory Data}

We use 467,911 GPS trajectories recorded between August~31 and September~30, 2018.
Each trajectory consists of timestamped $(latitude, longitude)$ points from mobile devices.
The raw data is processed through:
\begin{enumerate}
    \item \textbf{Cleaning}: Speed filtering ($v > 2$\,m/s), 5-minute gap splitting into separate trips.
    \item \textbf{Map matching}: Hidden Markov Model-based matching~\cite{newson2009hidden} using the Leuven map-matching library~\cite{meert2018hmm}, producing sequences of traversed link IDs with entry timestamps.
    \item \textbf{Temporal binning}: Counting traversals per link per 15-minute window, yielding a $2{,}880 \times 63{,}118$ sparse count matrix (2.06\% density).
\end{enumerate}

\subsection{The Sparsity Problem}
\label{sec:sparsity}

In any 15-minute window, only 1--3\% of network links have GPS observations.
The raw count vector $c_t$ for timeframe~$t$ has $c_t[i] = 0$ for $\sim$97\% of links.
Directly using $c_t / \|c_t\|_1$ as a traffic distribution assigns zero probability to most links, making the Mean Relative Error (MRE) undefined for unobserved links.

\subsection{Graph-Diffusion Smoothing}
\label{sec:smoothing}

We address sparsity through a one-step graph-diffusion filter~\cite{shuman2013emerging}:
\begin{equation}
    \boxed{c_{\text{smooth}} = \bigl[(1-\gamma)\,I + \gamma\,P_{\text{adj}}\bigr]\,\tilde{c}}
    \label{eq:smoothing}
\end{equation}
where $\tilde{c} = c + \alpha \cdot \mathbf{1}$ is the Laplace-smoothed count vector ($\alpha = 0.1$), $P_{\text{adj}}$ is the row-stochastic adjacency matrix, and $\gamma \in [0, 1]$ is the diffusion strength.

\textbf{Spectral interpretation.}
$P_{\text{adj}}$ has eigenvalues $\lambda \in [-1, 1]$.
The filter $(1-\gamma)I + \gamma P_{\text{adj}}$ has eigenvalues $(1-\gamma) + \gamma\lambda$.
The DC component ($\lambda = 1$) is preserved exactly; high-frequency noise ($\lambda \approx -1$) is attenuated by factor $|1-2\gamma|$.
At $\gamma = 0.26$, attenuation is 48\%.

\textbf{Cross-validation.}
We determine $\gamma$ via split-half cross-validation (49~random timeframes, seed~42), measuring both predictive MSE and downstream PageRank MRE.
The optimal $\gamma = 0.26$ minimises both metrics simultaneously, retaining 69\% of the original variance.
After smoothing, the matrix becomes $2{,}880 \times 99{,}716$ with 100\% density---every link has a positive probability estimate.


% ================================================================
\section{Two-Phase PageRank Model}
\label{sec:model}

\subsection{Standard PageRank}
The classical PageRank~\cite{page1999pagerank, brin1998anatomy} computes the stationary distribution of a random walk on a graph:
\begin{equation}
    v^{(k+1)} = d \cdot P^\top v^{(k)} + (1-d) \cdot E
\end{equation}
where $d$ is the damping factor, $P$ is a row-stochastic transition matrix, and $E$ is a teleportation vector (uniform in the original formulation, or empirical in our case).

\subsection{Travel-Time Self-Loops}
\label{sec:selfloops}

In a road network, vehicles spend time \emph{on} each link proportional to its traversal time $\ell_i / s_i$.
Standard PageRank treats all transitions as instantaneous, assigning equal weight to short connectors and long highways.
We introduce self-loops that model dwell time:
\begin{equation}
    \boxed{sw_i = \mu \cdot \frac{\ell_i}{s_i}}
    \label{eq:selfloop}
\end{equation}
where $\mu > 0$ is the dwell-time scaling parameter.
The transition probability from link~$i$ to successor~$j$ becomes:
\begin{equation}
    P[i, j] = \frac{w_{ij}}{sw_i + \sum_k w_{ik}}, \qquad
    P[i, i] = \frac{sw_i}{sw_i + \sum_k w_{ik}}
    \label{eq:transition}
\end{equation}

For a typical highway link ($\ell = 500$\,m, $s = 25$\,m/s, 3~successors):
\[
sw = 20 \times 500/25 = 400, \quad P[i,i] = \frac{400}{400+3} = 99.3\%
\]
After $k$ iterations, the retained mass is $P[i,i]^k$.
This naturally allocates more PageRank mass to links with longer traversal times---which are precisely the high-volume corridors that carry the most traffic.

\subsection{Two-Phase Block Structure}

We decompose the random walk into two phases:
\begin{itemize}
    \item \textbf{Up phase} (active driving): follows the directed adjacency with travel-time self-loops.
    \item \textbf{Down phase} (local dispersal): vehicles that have ``arrived'' disperse locally.
\end{itemize}
The $2N \times 2N$ transition matrix is:
\begin{equation}
    \boxed{\MTN = \begin{bmatrix} (1-\beta)\,P_{\text{up}} & \beta\,I_N \\ \mathbf{0} & P_{\text{down}} \end{bmatrix}}
    \label{eq:block}
\end{equation}
where $\beta \in (0,1)$ is the phase-drop rate.
At each step, a fraction $\beta$ of up-phase mass transitions to the down phase.
The optimal $\beta = 0.9$ implies the average trip has approximately $1/\beta \approx 1.1$ forward hops before dispersal.

\subsection{Power Iteration}

The full update rule is:
\begin{equation}
    \boxed{v^{(k+1)} = d \cdot \MTN^\top \, v^{(k)} + (1-d) \cdot \ETN}
    \label{eq:power}
\end{equation}
where $\ETN = [E_N; \,\mathbf{0}]$ injects the smoothed ground truth into the up phase only, and $d = 0.80$ is the damping factor.
Convergence is measured by $\|v^{(k+1)} - v^{(k)}\|_1 < 10^{-6}$ (typically 37--38 iterations).

The final distribution is obtained by collapsing phases:
\begin{equation}
    \vfinal = \frac{v_{\text{up}} + v_{\text{down}}}{\|v_{\text{up}} + v_{\text{down}}\|_1}
\end{equation}


% ================================================================
\section{Demand-Weighted Transitions}
\label{sec:demand}

\subsection{Motivation}
\label{sec:demand_motivation}

At $d = 0.80$, 20\% of the prediction comes from ground truth.
To reduce this dependence, we increase~$d$---but the topology-only matrix does not encode traffic demand patterns, so accuracy degrades rapidly (Table~\ref{tab:damping}).

A grid search over 576 parameter combinations ($\mu, \alpha_s, \alpha_l, \beta$) found \emph{zero feasible solutions} for $d > 0.90$ with MRE~$\leq 0.05$.

\subsection{The $\gamma$ Parameter}

We introduce a single parameter $\gamma$ that weights each transition by the destination link's average demand:
\begin{equation}
    \boxed{P(i \to j) \propto w_{ij} \cdot \text{demand}[j]^\gamma}
    \label{eq:gamma}
\end{equation}
where $\text{demand}[j] = \frac{1}{T}\sum_{t=1}^T E_N^{(t)}[j]$ is the time-averaged popularity of link~$j$ across a training set of $T$ timeframes.
This teaches the matrix that high-traffic links (e.g., I-15) should attract more flow than residential streets.

\subsection{Validation}

We validate through three independent train/test splits:
\begin{enumerate}
    \item \textbf{Temporal}: Train on Sept~1--15, test on Sept~16--30.
    \item \textbf{Day-of-week}: Train on Mon/Wed/Fri, test on Tue/Thu/Sat/Sun.
    \item \textbf{5-fold random CV}: 80\% train, 20\% test, repeated 5 times.
\end{enumerate}

Results (Table~\ref{tab:demand_validation}) confirm the improvement generalises: at $d = 0.96$, $\gamma = 0.20$, the average Overall~MRE across all splits is $\mathbf{0.0385}$---better than the original model at $d = 0.80$ (0.0466), while using only 4\% ground truth.

\begin{table}[t]
\centering
\caption{Overall MRE across three validation strategies.}
\label{tab:demand_validation}
\small
\begin{tabular}{lcccc}
\toprule
Configuration & Temporal & DoW & 5-Fold & \textbf{Avg} \\
\midrule
$d{=}0.80,\;\gamma{=}0$ & 0.0478 & 0.0457 & 0.0464 & 0.0466 \\
$d{=}0.91,\;\gamma{=}0.10$ & 0.0472 & 0.0452 & 0.0459 & \textbf{0.0461} \\
$d{=}0.94,\;\gamma{=}0.15$ & 0.0439 & 0.0420 & 0.0427 & \textbf{0.0429} \\
$d{=}0.96,\;\gamma{=}0.20$ & 0.0395 & 0.0377 & 0.0383 & \textbf{0.0385} \\
\bottomrule
\end{tabular}
\end{table}


% ================================================================
\section{Results}
\label{sec:results}

\subsection{Evaluation Protocol}

We evaluate on 49 random timeframes (seed~42) held out from parameter tuning.
The ground truth is the smoothed popularity vector $E_N$ for each timeframe.
The prediction is $\vfinal$ from the two-phase PageRank.
The primary metric is Mean Relative Error:
\begin{equation}
    \MRE = \frac{1}{|S|}\sum_{i \in S} \frac{|E_N[i] - \vfinal[i]|}{E_N[i] + \epsilon}
\end{equation}
computed over all links (Overall) or the top-100 busiest links (Top-100).

\subsection{Main Results}

\begin{table}[t]
\centering
\caption{Evaluation metrics (49 timeframes, $\mu{=}20$, $\beta{=}0.9$, $d{=}0.80$).}
\label{tab:main_results}
\small
\begin{tabular}{lrr}
\toprule
Metric & Mean $\pm$ Std & Range \\
\midrule
Top-100 MRE & $0.0312 \pm 0.0122$ & 0.020--0.048 \\
Overall MRE & $0.0464 \pm 0.0056$ & 0.041--0.057 \\
Pearson $r$ & $0.9969 \pm 0.0014$ & 0.995--0.998 \\
Rank overlap /100 & $94.4 \pm 4.5$ & 86--99 \\
\bottomrule
\end{tabular}
\end{table}

Table~\ref{tab:main_results} shows the model achieves Top-100~MRE of 3.1\%, Pearson correlation 0.997, and correctly identifies 94 out of 100 busiest links on average.

\subsection{Effect of Damping Factor}

\begin{table}[t]
\centering
\caption{Damping factor sweep (49 timeframes, $\gamma{=}0$).}
\label{tab:damping}
\small
\begin{tabular}{cccc}
\toprule
$d$ & Top-100 MRE & Overall MRE & Pearson $r$ \\
\midrule
0.80 & 0.0312 & 0.0464 & 0.9969 \\
0.84 & 0.0409 & 0.0605 & 0.9949 \\
0.88 & 0.0558 & 0.0817 & 0.9912 \\
0.92 & 0.0820 & 0.1177 & 0.9829 \\
0.96 & 0.1443 & 0.1947 & 0.9565 \\
\bottomrule
\end{tabular}
\end{table}

Without demand weighting ($\gamma = 0$), MRE increases monotonically with $d$ (Table~\ref{tab:damping}).
With $\gamma = 0.20$, the model achieves Overall~MRE~$= 0.0383$ at $d = 0.96$ (Table~\ref{tab:demand_validation})---a 5$\times$ reduction in ground-truth dependence with improved accuracy.


% ================================================================
\section{Event Detection and Prediction}
\label{sec:events}

We validate the model against two documented real-world events:
an unplanned traffic crash and a planned sporting event.

\subsection{I-15 Multi-Vehicle Crash (September~20, 2018)}
\label{sec:crash}

A four-vehicle collision blocked multiple southbound lanes of I-15 at 14400~South, Draper, from approximately 7:30--11:00~AM~\cite{gephardtdaily2018}.
We compare crash-day PageRank mass on the 21~I-15~SB links against three baseline Thursdays (Sept~6, 13, 27).

\begin{table}[t]
\centering
\caption{I-15 crash: PageRank deviation from baseline.}
\label{tab:crash}
\small
\begin{tabular}{llrrr}
\toprule
Time & Phase & I-15 SB & I-15 NB & Surface \\
\midrule
06:00 & Pre-crash & $+1.7\%$ & $+96.5\%$ & $+10.4\%$ \\
08:30 & Peak & $\mathbf{-18.9\%}$ & $\mathbf{-92.7\%}$ & $-58.6\%$ \\
10:30 & Rebound & $\mathbf{+291.4\%}$ & $-40.9\%$ & $-8.3\%$ \\
11:30 & Recovery & $+36.7\%$ & $+37.1\%$ & $+15.2\%$ \\
\bottomrule
\end{tabular}
\end{table}

The model captures the complete crash lifecycle (Table~\ref{tab:crash}):
pre-crash calibration ($+1.7\%$, confirming fair comparison), peak disruption ($-18.9\%$~SB, $-92.7\%$~NB from rubbernecking), clearance rebound ($+291.4\%$~SB as pent-up traffic releases), and recovery ($+36.7\%$).
The demand-weighted model ($d = 0.96$, $\gamma = 0.20$) produces virtually identical deviations ($-18.9\%$, $+292.4\%$), confirming crash detection with only 4\% ground truth.

\subsection{Utah vs.\ \#10 Washington Football Game (September~15, 2018)}
\label{sec:game}

A near-capacity crowd of 47,445 attended the Saturday night game at Rice-Eccles Stadium (kickoff 8:00~PM)~\cite{utahutes2018}.
Documented road closures include 500~South lane reduction and Guardsman~Way full closure~\cite{utahathletics2018}.
We compare game day against four baseline Saturdays (Sept~1, 8, 22, 29) on 93~event-affected links.

The model detects the full game-day cycle:
$-85.9\%$ on 500~South at 14:00 (road closures),
$+463.4\%$ at 19:00 (arrival surge),
$+247.9\%$ at stadium core at 23:00 (post-game departure).


% ================================================================
\section{The Role of $\mu$: Per-Link Disruption Control}
\label{sec:mu}

\subsection{Why $\mu$ Is the Only Effective Control Parameter}

The self-loop probability $P[i,i] = sw_i / (sw_i + \sum_k w_{ik})$ is dominated by $sw_i = \mu \cdot \ell_i / s_i$ at $\mu = 20$ (typically $> 99\%$).
We tested five types of intervention on crash-affected links:

\begin{table}[t]
\centering
\caption{Effect of different interventions on I-15~SB corridor mass.}
\label{tab:interventions}
\small
\begin{tabular}{lrc}
\toprule
Intervention & What changes & Max effect \\
\midrule
Speed $= 0.1$\,m/s & Outgoing weights & $-0.44\%$ \\
Lanes $= 1$ & Outgoing weights & $-0.54\%$ \\
Remove from adjacency & Zero outgoing & $+0.67\%$ \\
Zero demand weight & Zero attraction & $+0.67\%$ \\
\textbf{$\mu = 0.01$} & \textbf{Self-loop: $99.5\% \to 9.6\%$} & $\mathbf{-64.4\%}$ \\
\bottomrule
\end{tabular}
\end{table}

Only $\mu$ produces meaningful changes (Table~\ref{tab:interventions}) because it operates on the 99.5\% retained mass, while all other parameters affect only the 0.5\% outgoing flow.

\subsection{Crash Prediction Without Crash-Day Data}

Using September~13 (normal Thursday) as input and reducing $\mu$ on the 21~I-15~SB links:

\begin{table}[t]
\centering
\caption{I-15~SB prediction error at key crash times.}
\label{tab:mu_crash}
\small
\begin{tabular}{lccrr}
\toprule
Time & $\mu{=}20$ & Best $\mu$ & Best err & Improv. \\
\midrule
08:00 & 34.1\% & 0.10 & \textbf{0.0\%} & $+100\%$ \\
08:30 & 179.1\% & 0.01 & \textbf{11.5\%} & $+93.6\%$ \\
08:45 & 380.7\% & 0.00 & \textbf{26.4\%} & $+93.1\%$ \\
10:15 & 142.2\% & 0.01 & \textbf{2.6\%} & $+98.1\%$ \\
\bottomrule
\end{tabular}
\end{table}

During the active diversion phase, $\mu$-control reduces corridor error by 91--100\% (Table~\ref{tab:mu_crash}).
The same approach on the football game yields 94--98\% error reduction on stadium-area links during closure phases.

\subsection{Cross-Validated Baseline: Tuesday $\mu$ Profile}

To confirm that $\mu = 20$ is genuinely optimal on normal days, we perform leave-one-out cross-validation across four Tuesdays (no known events).
For each of 96~timeframes, we sweep $\mu$ and find the value minimising Overall~MRE.

\textbf{Result}: The optimal $\mu$ stays in the range 10--20 across all 96~slots, with only 1\% improvement over fixed $\mu = 20$.
The correlation between raw ground-truth gap (between Tuesdays) and model prediction error is $r = 0.9999$---the model adds less than 4\% distortion to the input data.
The model error is 96.6\% of the raw input gap, confirming that the error comes from week-to-week GPS variation, not model miscalibration.

This establishes the baseline: $\mu \approx 20$ on normal days, so any significant deviation signals a real disruption.

\subsection{Practical Protocol}

For real-time event detection and scenario analysis:
\begin{enumerate}
    \item Build the demand-weighted matrix ($\gamma = 0.20$) once using historical data.
    \item For \textbf{detection}: run PageRank with current GPS data and compare to baseline. Deviations beyond $\pm 2\sigma$ flag anomalies.
    \item For \textbf{prediction}: when a disruption is reported on specific links, set $\mu \to 0.01$ on those links and rerun PageRank with the most recent normal-day data. The output predicts the new traffic equilibrium in $<$10\,ms on GPU.
    \item \textbf{Severity mapping}: $\mu = 20$ (normal), $\mu \sim 1$ (partial restriction), $\mu \sim 0.01$ (full closure).
\end{enumerate}


% ================================================================
\section{Conclusion}
\label{sec:conclusion}

We presented a Two-Phase PageRank model with travel-time self-loops and demand-weighted transitions for urban traffic distribution estimation.
The model achieves 3.1\% Top-100~MRE on 49 held-out timeframes and detects real-world disruptions (crashes, sporting events) through PageRank deviation analysis.

Three parameters control the model:
\begin{itemize}
    \item \textbf{Damping $d$}: balances matrix topology against ground truth (optimal: 0.96 with demand weighting).
    \item \textbf{Demand sensitivity $\gamma$}: encodes structural traffic demand into the matrix (optimal: 0.20, validated across 3~independent splits).
    \item \textbf{Dwell-time $\mu$}: controls per-link mass retention---the only effective parameter for local disruption modelling (default: 20, disrupted: 0.01).
\end{itemize}

The key insight is the separation of concerns: $\gamma$ handles global prediction accuracy, while $\mu$ handles local disruption response.
Together, they enable a model that uses only 4\% ground truth yet achieves high accuracy and can predict event-driven traffic redistribution without any event-day data.

Future work includes time-varying $\mu$ schedules matched to incident timelines, extension to multi-modal networks, and integration with real-time incident detection systems.

\bibliographystyle{plainnat}
\bibliography{references}

\end{document}
